\section{评估证据}\label{chapter:eval}

\subsection{问题与证据类别}\label{sec:eval:sample}

评估区分四个问题。\emph{合法性}询问非法制度变更能否被拒绝、合法变更
能否原子提交。\emph{活性}询问提议、攻击、裁决、转换、修复和重放能否
推进。\emph{检测性能}询问已知故障的检出率以及正常样例的误拒率。
\emph{有效性}询问受治理的协同演化能否改善研究质量、校准或成本。试验
主要回答合法性,固定故障样例只演练特定活性路径。检测性能和有效性需要
有正确答案的样例、独立运行及对照条件。

\begin{table}[t]
  \centering
  \small
  \begin{tabularx}{\columnwidth}{@{}lXX@{}}
    \toprule
    \textbf{证据} & \textbf{支持的主张} & \textbf{不支持的主张} \\
    \midrule
    1,124 项相关试验 & 已实现规则在受测输入上的行为 &
      未测试输入上的安全性或研究质量 \\
    五个故障、一个正常样例 & 五条指定路径触发且正常样例通过 &
      误检率、漏检率或未知故障覆盖率 \\
    探索性运行记录 & 部分长路径能够执行 &
      可复现性、总体活性或优越性 \\
    尚未运行的对照研究 & 无 &
      质量、校准、成本或组件级因果效应 \\
    \bottomrule
  \end{tabularx}
  \caption{当前证据支持和不支持的主张。}
  \label{tab:eval-status}
\end{table}

\subsection{生命周期演示}\label{sec:eval:construction}

成果说明中记录的精确选择条件收集了直接涉及 RQGM 与最终论文闸门的
1,124 项试验；本工作树中全部通过。完整套件收集 4,805 项：4,787 项通过，
16 项跳过，2 项为预期失败。聚焦集合无重叠地分为：最终论文验证 64；
内核、转换、纪元、状态、创始和模式 285；对抗与治理 126；修复与净室
66；效用演化 85；评估样例与指标 157；提示词演化 93；论文角色 94；
上下文、预算、元演化、提案与适配器 154，共 1,124 项。机器可读成果物
列出每个试验标识。
\Cref{sec:appendix:specification} 将机制对应到组件生命周期表、权限矩阵
与恢复过程。

此前引用的完整 SHA 无法从公开仓库解析，也不标识这个未提交修订，故撤回
公开成果物主张。本工作树现包含本地内容寻址复现包：源文件摘要、依赖清单、
精确选择条件、全部收集标识、完全展开的机器可读权限矩阵及聚焦/完整运行
日志。脏工作树由聚合内容摘要
标识，而不假装当前 Git SHA 指向它。但它不是不可变公共标签、DOI 归档、
容器或 CI 记录，第三方尚不能取得。论文引用公共成果标识前，必须与完整
提示词材料一起永久发布和归档。

两个预期失败是独立的架构债务：一个命令层模块仍导入可视化层，评估器仍
绕过共享模型客户端直接调用提供者库。它们不绕过 RQGM 内核检查，但保留
在完整日志中且不计为通过。

试验数量是实现规模和防止回归的证据,不是科学性能。当前环境没有覆盖率
测量工具,因此代码行覆盖率、分支覆盖率和变异试验得分未知。总数也不能
证明每项不变量至少有一项试验。本文不把总试验数作为安全性的替代指标。

初始注册数取决于配置。基础研究配置注册 31 个提示词、一个效用规则和
20 个组件。启用论文存档后,自我偏好检查、论文审稿和论文写作各增加一个
提示词和组件,总计为 35 个提示词或规则记录、23 个组件。基础集合包含
已注册研究生成者；此前“32 与 19”是修正前计数。

较早的探索性执行观察到边界事务、效用规则替换、从保存的原始评价轴
重新评分,以及论文角色后继者的采用。但是,本报告没有附带完整保存点、
精确的服务端模型版本、解码设置、随机种子、独立运行次数、任务说明和
成本。因此这些只是实现路径观察,不是可复现的实验结果。

\subsection{故障注入与构造保证}\label{sec:eval:repro}

固定且有正确答案的样例覆盖五条故障路径:无法从运行值重新计算的指标
操纵、缺少证据的过度主张、包含封存内容的提示词候选、隔离再生成的
禁止输入,以及依赖已退役提示词的记录混入候选集合。它们分别通过预期的
数值证据、候选、内核或影响排除路径被拒绝。一个正常数值主张没有被
错误拒绝。

额外完整性负例修改事件类型、标识、事务标识、载荷与前驱摘要，重排已提交
事件，并制造事务成员不匹配；这些情况全部被拒绝。T16 试验还验证强制关闭/
开启纪元会产生新指纹，机器可读权限导出也与稀疏允许表逐项核对。这些是
确定性回归用例，不是事件模糊测试或模型检查。

该结果只表明这些固定样例上 5/5 个故障被检出且 1/1 个正常样例通过。
样本过小,每种故障只有一个难度,所以本文不报告误检率或漏检率估计。
虚构先行研究、干扰裁决路径、责任对象错误绑定,以及保留语义的隔离
再生成污染尚无完整的带标签检测实验。

构造保证也有边界。固定代码会拒绝转换器以外的注册表变更、表中没有的
转换、候选权限扩张、未裁决攻击导致的扣分,以及未提交事务的重放。这是
针对已保存输入的内核确定性,不是模型输出的再生成性。版本二审计链绑定
所有影响重放语义的信封字段并检测已测试的中间修改,但没有外部固定点时,
无法检测整个保存点回退到旧的合法前缀。T16 现在强制关闭当前纪元，并在
同一事务中开启具有新指纹的纪元。

\subsection{局限}\label{sec:eval:limits}

达到可接收水平的研究必须保存预注册任务集、精确模型及服务端版本、
温度、采样方法、最大输出量、工具版本、随机种子、独立运行次数、纪元和
节点数、候选数、攻击/动议/裁决/退役数,以及令牌、时间和金钱成本。每项
比率都需给出分母、运行间方差和置信区间。

我们实现了与原 RQGM 研究对齐的五个论文比较条件。P0 仅进化写作角色并
固定审稿角色；P1 只增加审稿角色替换；P2 共同进化两个角色，但保留被替换
审稿角色产生的效用记录；P3 将这些旧效用排除出选择，构成完整 RQGM；
P4 在 P3 上增加具有阻断能力的立宪内核。五个条件使用相同的论文归档、
预算、随机种子和模型设置，只是评估预设，并非新的生产运行模式。主比较
为 P0/P3/P4，P1/P2 用于机制消融。论文写作和固定的事后 rubric 评审可使用
分别钉定的模型；计划中的实验以 Claude Code 写作，并以 Codex 执行独立于
协同进化审稿角色的固定面板评审。本文尚未包含真实模型的多随机种子运行，
因此这是已实现的评估协议，而不是实验结果。
这些条件在机制上遵循原研究，但随附预算仅为八个纪元、每个纪元最多
12 次扩展，并未复现原论文的 12,288 次评估规模
\cite{iacob2026red}。投稿实验必须预注册共同的评估调用预算，并报告
各条件实际产生的调用次数。

进一步的组件消融仍需逐一去除历史样例重放、责任绑定、隔离再生成、可变
效用规则、基准资料和固定稿件闸门。责任归属和影响修复还需要带标签的合成
因果图，用于测量目标准确率、影响范围的精确率与召回率，以及过度删除。

因此,现阶段可成立的结论很有限:五个固定故障样例触发了预期检查,
事件完整性负例被拒绝，1,124 项相关试验全部通过。研究质量提升、一般检测性能、长期活性和
成本效益均尚未得到证明。
